\documentclass[tikz,border=3mm]{standalone}
% Spolecny styl vlastnich vektorovych ilustraci; preklad jednotlivych stackup_*.tex.
\usepackage[czech]{babel}
\usepackage[utf8]{inputenc}
\usepackage[T1]{fontenc}
\usepackage{tgheros}
\renewcommand{\familydefault}{\sfdefault}
\definecolor{sig}{RGB}{244,185,75}
\definecolor{gnd}{RGB}{107,183,141}
\definecolor{pwr}{RGB}{228,129,112}
\definecolor{frfour}{RGB}{245,237,214}
\tikzset{copper/.style={draw=black!55,line width=0.4pt},
 layerlabel/.style={anchor=east,font=\small},
 note/.style={font=\small,align=center}}

\usepackage{amsmath}
\usetikzlibrary{arrows.meta,decorations.pathmorphing}
\begin{document}
\begin{tikzpicture}[font=\small,>=Latex]
 \draw[thick] (0,3)--(3,3) (0,0)--(8.5,0);
 \node[anchor=east] at (0,3) {VIN};
 \node[anchor=east] at (0,0) {PGND};
 % Vstupni a vystupni kondenzator.
 \foreach \x in {.6,7.5} {
  \draw[thick] (\x,3)--(\x,1.65) (\x,1.35)--(\x,0);
  \draw[thick] (\x-.3,1.65)--(\x+.3,1.65) (\x-.3,1.35)--(\x+.3,1.35);
 }
 \node[anchor=east] at (.2,1.5) {$C_{\text{IN}}$};
 \node[anchor=west] at (7.9,1.5) {$C_{\text{OUT}}$};
 % Funkcni bloky spinacu, ridici vyvody nejsou zakreslene.
 \draw[thick] (3,3)--(3,2.75) (3,.75)--(3,0);
 \filldraw[fill=black!6,thick] (2.55,2) rectangle (3.45,2.75);
 \filldraw[fill=black!6,thick] (2.55,.75) rectangle (3.45,1.5);
 \node at (3,2.38) {QH}; \node at (3,1.12) {QL};
 \draw[red!75!black,line width=3pt] (3,2)--(3,1.5) (3,1.75)--(5,1.75);
 \fill[red!75!black] (3,1.75) circle (.07);
 \node[above,text=red!75!black] at (4.2,1.85) {SW};
 \draw[thick,decorate,decoration={coil,aspect=.5,segment length=5pt,amplitude=4pt}] (5,1.75)--(6.2,1.75);
 \node[above] at (5.6,2.05) {L};
 \draw[thick] (6.2,1.75)--(7.5,1.75) (7.5,1.75)--(7.5,3)--(8.5,3);
 \node[anchor=west] at (8.5,3) {VOUT};
 \node[anchor=west,font=\scriptsize] at (0,-.55) {Červeně: spínací uzel SW.};
\end{tikzpicture}
\end{document}
